\documentclass[11pt,a4paper]{article}

\usepackage[margin=2.3cm]{geometry}
\usepackage[T1]{fontenc}
\usepackage{lmodern}
\usepackage[utf8]{inputenc}
\usepackage{microtype}
\usepackage{xcolor}
\usepackage{enumitem}
\usepackage{booktabs}
\usepackage{longtable}
\usepackage{array}
\usepackage{titlesec}
\usepackage{fancyhdr}
\usepackage[colorlinks=true,linkcolor=accent,urlcolor=accent,citecolor=accent]{hyperref}

\definecolor{accent}{RGB}{20,70,120}
\definecolor{soft}{RGB}{240,244,248}

\setlist[itemize]{leftmargin=1.3em,itemsep=2pt,topsep=3pt,parsep=0pt}
\setlist[enumerate]{leftmargin=1.5em,itemsep=2pt,topsep=3pt,parsep=0pt}
\setlength{\parindent}{0pt}
\setlength{\parskip}{5pt}
\renewcommand{\arraystretch}{1.25}

\titleformat{\section}{\Large\bfseries\color{accent}}{\thesection}{0.6em}{}
\titleformat{\subsection}{\large\bfseries\color{accent}}{\thesubsection}{0.6em}{}
\titlespacing*{\section}{0pt}{16pt}{6pt}
\titlespacing*{\subsection}{0pt}{12pt}{4pt}

\pagestyle{fancy}
\fancyhf{}
\fancyhead[L]{\small Forecasting Programme Plan}
\fancyhead[R]{\small Alqueva PSP--PV--BESS Trading Optimizer}
\fancyfoot[C]{\small\thepage}
\renewcommand{\headrulewidth}{0.3pt}

\newcommand{\exitcrit}[1]{\par\smallskip\noindent\colorbox{soft}{\parbox{\dimexpr\linewidth-2\fboxsep}{\textbf{Exit criterion.} #1}}\par}

\begin{document}

% ---------------------------------------------------------------- title
\begin{center}
{\LARGE\bfseries\color{accent} A Professional Forecasting Programme\\[2pt]for an Electricity-Market Company}\\[8pt]
{\large From raw data to monitored production forecasts}\\[6pt]
{\normalsize Project: Alqueva PSP--PV--BESS 24-hour energy trading optimizer\\
Document status: plan for review \quad|\quad Date: 10 October 2026}
\end{center}

\vspace{4pt}
\hrule
\vspace{6pt}

\section{Purpose and how to read this plan}

This document describes, step by step, how a professional energy-market team would build,
validate, deploy and maintain machine-learning forecasts, and how we will apply the same
process to the forecasters of this project. It is deliberately complete: no step is
shortened, because the work can be done on a computer without a deadline.

The plan has six parts (A to F) and eighteen phases (0 to 17). Every phase lists its steps and an
\emph{exit criterion}: the evidence that must exist before the next phase starts. We complete the phases
in order, one at a time.

\subsection*{What the literature and industry practice say}

\begin{itemize}
\item \textbf{Benchmark honestly.} The standard review of day-ahead price forecasting
  (Lago et al., 2021) identifies the weak points of the field: short or single-market test
  samples, no comparison with simple benchmarks, unsuitable accuracy metrics and no test of
  statistical significance.
\item \textbf{Backtest the way production will run.} Validation always trains on the past and tests on the
  future with rolling origins; shuffled cross-validation leaks information. A model that cannot beat a
  naive benchmark has not found structure in the data.
\item \textbf{Probabilistic forecasts need proper scoring.} Quantile forecasts are scored with the
  pinball loss; the CRPS is approximated by averaging pinball losses over a quantile grid; coverage is
  tested separately; models are compared with Diebold--Mariano tests. Decision-based metrics
  (profit) are added because a better CRPS does not always mean better decisions.
\item \textbf{Balancing markets are harder.} Studies of German aFRR and imbalance prices find no
  universally best model and poor results for activated volumes; Nordic mFRR studies find that
  extra model or feature complexity does not reliably help. We found no Iberian study, so our own
  evidence is needed.
\item \textbf{Production needs MLOps.} Versioned features, a model registry, data validation, drift
  monitoring and champion--challenger promotion are standard in energy-trading forecasting systems.
\end{itemize}

\subsection*{Where the project stands today}

The following results were measured earlier on real market data (last twelve months as the test window).

\begin{center}
\begin{tabular}{lp{7.6cm}}
\toprule
\textbf{Forecast} & \textbf{Result versus the naive benchmark}\\
\midrule
Day-ahead price & about 2.6\% better\\
IDA1--3 and XBID & 1\% to 5\% worse\\
aFRR capacity price & very poor (no one-day lag feature)\\
mFRR price & about 0.2\% better\\
PV (irradiance and temperature) & worse than naive\\
Reservoir inflow & worse than naive\\
\bottomrule
\end{tabular}
\end{center}

Existing foundations: real data rebuilt from OMIE and REN, leakage fixes, walk-forward cross-validation,
naive baselines, autoregressive evaluation of the intraday forecasters. Missing: probabilistic
forecasts, fundamental features (load, wind, solar, fuel prices), significance tests, experiment
tracking, drift monitoring, a forecast archive and decision-value evaluation.

% ================================================================ PART A
\section{Part A --- Foundations}

\subsection{Phase 0. Charter and decision mapping}
\begin{itemize}
\item List every decision that consumes a forecast: day-ahead bid, IDA1--3 and XBID trades, aFRR and mFRR
  capacity offers, PV and inflow planning, risk limits.
\item For each forecast define the \textbf{as-of time}: exactly which information exists at the moment of the
  decision (day-ahead gate 12:00 CET, IDA1 15:00 CET, IDA2 22:00 CET, IDA3 10:00 CET, reserve gates, and so on).
  This is the main protection against leakage.
\item Define horizon, resolution (hourly and 15-minute), output format (point, quantiles or scenarios) and the
  consumers of each forecast.
\item Define success in euros (value of the forecast in the optimiser), not only in MAE.
\item Fix the forecast catalogue: targets, owners, update schedule, fallback.
\end{itemize}
\exitcrit{A signed-off forecast catalogue with an as-of time and a consumer for every target.}

\subsection{Phase 1. Data inventory and access}
\begin{itemize}
\item For every source (OMIE prices, REN reserve data, weather observations and forecasts, TSO fundamentals,
  fuel and carbon prices, reservoir data): history length, update delay, licence, API limits, cost.
\item Record \textbf{vintages}: the forecast that was issued at the time, not only the final actual value
  (point-in-time data), for load, wind, solar and weather forecasts.
\item Build a data dictionary (variable, unit, resolution, time zone, source, definition) and a lineage map.
\item Identify gaps and substitutes, and the legal use of each dataset.
\end{itemize}
\exitcrit{Every input has an owner, a timestamp rule, a licence note and a quality note.}

\subsection{Phase 2. Ingestion and storage}
\begin{itemize}
\item Three zones: \emph{raw} (immutable, as downloaded), \emph{cleaned}, and \emph{feature-ready}, each versioned.
\item Incremental, repeatable loads: retries, idempotent writes, checksums, load logs.
\item Time handled properly: UTC internally; 23- and 25-hour days; 92- and 100-period quarter-hour days;
  the move from hourly to 15-minute products (2025); the 2024 intraday-auction reform; the 2026 aFRR platform
  change. Each is stored as a \textbf{regime marker}.
\item Data contracts: schema, types, ranges and freshness rules checked on every load.
\end{itemize}
\exitcrit{Any past date can be rebuilt exactly from the raw zone.}

\subsection{Phase 3. Data quality and cleaning}
\begin{itemize}
\item Automated checks: schema, ranges, duplicates, gaps, stale feeds, unit consistency, DST artefacts.
\item Price spikes and negative prices are information: flag them, never delete them blindly.
\item Written imputation rules per variable that use only information available at the as-of time.
\item Keep the provenance label (live, synthetic) on every row; synthetic rows are excluded from training and testing.
\item A daily data-quality score and an automatic report.
\end{itemize}
\exitcrit{A quality report and a pass/fail gate that must be green before any model run.}

\subsection{Phase 4. Exploratory analysis and market understanding}
\begin{itemize}
\item Seasonality (daily, weekly, annual), autocorrelation, stationarity, structural breaks (the 2022 energy
  crisis, solar price cannibalisation, the April 2025 blackout effects).
\item Relationships between markets: day-ahead, IDA1--3, XBID, reserves, imbalance.
\item Frequency of spikes and negative prices, distribution shape, heteroscedasticity.
\item Effective sample size after the regime changes; which history is still representative.
\end{itemize}
\exitcrit{A written findings note that justifies the modelling choices of the next parts.}

% ================================================================ PART B
\section{Part B --- Measurement before modelling}

\subsection{Phase 5. Problem formulation per target}
\begin{itemize}
\item Output type: point, quantile or scenario; direct versus recursive multi-step; per-hour models versus a
  full 24-hour vector.
\item Targets as price levels or as spreads to the day-ahead price (important for intraday).
\item A variance-stabilising transform that allows negative prices (for example the inverse hyperbolic sine).
\item Consistency between hourly and 15-minute forecasts.
\item The information set per target, written down and tested.
\end{itemize}

\subsection{Phase 6. Evaluation framework (the contract, built first)}
\begin{itemize}
\item Rolling-origin backtests with a realistic retraining schedule and a time gap between training and test.
\item A \textbf{locked final test period} that nobody tunes on.
\item Point metrics: MAE, RMSE and scaled errors against the naive forecast (percentage errors are avoided because
  prices can be near zero or negative).
\item Probabilistic metrics: pinball loss on a stated quantile grid, CRPS, coverage tests, interval scores,
  PIT histograms.
\item Significance: Diebold--Mariano tests with a HAC variance, corrected for testing many hours and models.
\item Benchmarks per target: same hour yesterday, same hour last week, a simple regression (ARX), then the candidates.
\item \textbf{Decision value}: profit of the Alqueva optimiser when fed each forecast (naive, model, perfect foresight).
\item Automated leakage tests and reproducibility tests.
\end{itemize}
\exitcrit{One command scores any forecaster in the same way and produces the same report.}

% ================================================================ PART C
\section{Part C --- Modelling}

\subsection{Phase 7. Feature engineering and feature store}
\begin{itemize}
\item Calendar and holidays for Portugal and Spain, bridge days, DST.
\item Lags and rolling statistics that respect the as-of time.
\item Fundamentals: load, wind and solar forecasts, residual load, gas and carbon prices, interconnection, hydro
  levels, weather forecasts.
\item A feature catalogue with each feature's availability time, and point-in-time correctness tests.
\item Feature selection with stability checks across periods; versioned feature definitions.
\end{itemize}

\subsection{Phase 8. Baselines and candidate models}
\begin{itemize}
\item Statistical: seasonal naive, ARX, LASSO-estimated ARX with several calibration windows (a strong benchmark in
  the review), generalised additive models.
\item Machine learning: gradient boosting (LightGBM, XGBoost), random forest, quantile boosting.
\item Deep learning: feed-forward networks, LSTM and GRU, N-BEATS, Temporal Fusion Transformer, distributional networks.
\item Foundation models used zero-shot, as a benchmark only.
\item Rule: a complex model must beat the simple one significantly, otherwise it is dropped.
\end{itemize}

\subsection{Phase 9. Training, tuning and ensembling}
\begin{itemize}
\item Hyper-parameters tuned inside the rolling backtest (nested, time-aware), never on the test period.
\item Forecast combinations: simple averages often win and are tested first.
\item Comparison of calibration-window lengths and retraining frequencies.
\end{itemize}

\subsection{Phase 10. Probabilistic forecasts and scenarios}
\begin{itemize}
\item Quantile models, quantile regression averaging, distributional networks.
\item Conformal prediction for reliable intervals (recent studies apply it to day-ahead and balancing prices).
\item Scenario generation that keeps the correlation across hours, feeding the stochastic CVaR bidding model.
\item Calibration diagnostics and automatic recalibration.
\end{itemize}
\exitcrit{Intervals reach their nominal coverage out of sample, and scenarios pass the correlation checks.}

% ================================================================ PART D
\section{Part D --- Target-specific work}

\subsection{Phase 11. One workstream per target}
\begin{longtable}{>{\raggedright\arraybackslash}p{3.2cm}>{\raggedright\arraybackslash}p{12.2cm}}
\toprule
\textbf{Target} & \textbf{Focus}\\
\midrule
\endhead
Day-ahead price & Biggest revenue lever; the full pipeline is built here first (Phases 5--10).\\
IDA1--3 and XBID & Model the spread to the day-ahead price; use the outcome of the earlier gate as information; market-microstructure features.\\
aFRR and mFRR capacity prices & Day-before lag, merit-order features, bid and award behaviour; handle the 2026 product and platform change as a regime; expect modest gains.\\
Imbalance and activation & Probabilistic, low expectations; value is mainly in risk estimates.\\
PV & Forecast irradiance and temperature with weather-model post-processing, then apply the physics model; compare with satellite irradiance; probabilistic.\\
Reservoir inflow & Hydrological view (rainfall, runoff, irrigation rules); find out why the current forecast is worse than naive.\\
\bottomrule
\end{longtable}

% ================================================================ PART E
\section{Part E --- Industrialisation}

\subsection{Phase 12. Experiment tracking and reproducibility}
\begin{itemize}
\item Every run records code version, data snapshot, configuration, random seed, metrics and plots.
\item Automatic model cards (purpose, data, metrics, limits).
\end{itemize}

\subsection{Phase 13. Model selection, governance and acceptance gates}
\begin{itemize}
\item A model is promoted only if it beats the benchmark with statistical significance \emph{and} adds decision value.
\item Champion--challenger rule, documented assumptions and a written fallback.
\item Review and sign-off record for each promotion.
\end{itemize}

\subsection{Phase 14. Production pipeline}
\begin{itemize}
\item Scheduling tied to each market gate, retries, data-freshness limits, automatic fallback to a simple forecast
  when data are late.
\item A \textbf{forecast archive} with timestamps (essential for later evaluation).
\item Tests in continuous integration, containers, a fixed validated interface to the optimiser.
\end{itemize}

\subsection{Phase 15. Monitoring and maintenance}
\begin{itemize}
\item Input drift (population stability index, Kolmogorov--Smirnov), error drift, bias and calibration drift
  computed on realised data.
\item Alert thresholds, a retraining policy (scheduled and triggered), shadow runs of challengers.
\item Incident runbook and a periodic review.
\end{itemize}

% ================================================================ PART F
\section{Part F --- Value and communication}

\subsection{Phase 16. Decision value}
\begin{itemize}
\item Profit with perfect foresight versus model versus naive forecasts in the optimiser.
\item Sensitivity of profit to forecast error and the value in euros of each improvement.
\item Stochastic bidding with quantile forecasts versus point forecasts.
\end{itemize}

\subsection{Phase 17. Documentation and reporting}
\begin{itemize}
\item Methodology note, model cards, and the existing dashboard page extended with these metrics.
\item A short narrative of the work for technical interviews.
\end{itemize}

% ================================================================ order
\section{Proposed order of work}
\begin{enumerate}
\item Phases 0 and 1, plus an \textbf{audit} of today's forecasters (no new modelling code).
\item Phases 2 to 4: data layer, data quality, exploratory analysis.
\item Phase 6: the evaluation framework, before any new model.
\item The day-ahead price end to end (Phases 5 and 7 to 10), then reserves, then intraday, PV and inflow.
\item Phases 12 to 15, then Phase 16 to show the value.
\end{enumerate}

\section{Decisions needed before starting}
\begin{enumerate}
\item \textbf{Order of targets.} Proposed: day-ahead price, then aFRR and mFRR, then intraday, PV and inflow.
\item \textbf{Data access.} Are Spanish and Portuguese TSO data (load, wind and solar forecasts) and a weather
  forecast source available, or only freely available data?
\item \textbf{Tools.} Open-source Python (pandas, scikit-learn, LightGBM, Optuna, MLflow-style tracking) or specific tools?
\item \textbf{Tracking.} This document is the master checklist; each completed phase is recorded with its exit evidence.
\end{enumerate}

% ================================================================ references
\section*{References}
\begin{itemize}
\item J.~Lago, G.~Marcjasz, B.~De~Schutter, R.~Weron, \emph{Forecasting day-ahead electricity prices: A review of
  state-of-the-art algorithms, best practices and an open-access benchmark}, Applied Energy 293 (2021).
  \url{https://arxiv.org/abs/2008.08004}
\item J.~Nowotarski, R.~Weron, \emph{Recent advances in electricity price forecasting: A review of probabilistic
  forecasting}, Renewable and Sustainable Energy Reviews 81 (2018) 1548--1568.
\item \emph{The evolution of probabilistic price forecasting techniques: a review of the day-ahead, intra-day and
  balancing markets}. \url{https://arxiv.org/pdf/2511.05523}
\item \emph{Conformal prediction for electricity price forecasting in the day-ahead and real-time balancing market}.
  \url{https://arxiv.org/pdf/2502.04935}
\item \emph{Probabilistic forecasting of German electricity imbalance prices}. \url{https://arxiv.org/pdf/2205.11439}
\item \emph{Forecast evaluation for data scientists: common pitfalls and best practices}.
  \url{https://arxiv.org/pdf/2203.10716}
\item \emph{Greykite: deploying flexible forecasting at scale at LinkedIn} (rolling-origin backtesting).
  \url{https://arxiv.org/pdf/2207.07788}
\end{itemize}

\end{document}
